\documentclass[conference]{IEEEtran}
\IEEEoverridecommandlockouts

\usepackage{cite}
\usepackage{amsmath,amssymb,amsfonts}
\usepackage{algorithm}
\usepackage{algorithmic}
\usepackage{graphicx}
\usepackage{textcomp}
\usepackage{booktabs}
\usepackage{multirow}
\usepackage{stfloats}
\usepackage{flafter}
\usepackage{url}
\usepackage[caption=false,font=footnotesize]{subfig}
\usepackage[acronym,nonumberlist,nopostdot,nogroupskip]{glossaries}

\graphicspath{{../out/}}

\renewcommand{\topfraction}{0.92}
\renewcommand{\bottomfraction}{0.75}
\renewcommand{\textfraction}{0.06}
\renewcommand{\floatpagefraction}{0.70}
\setcounter{topnumber}{3}
\setcounter{bottomnumber}{2}
\setcounter{totalnumber}{5}

\setacronymstyle{long-short}
\newacronym{cu}{CU}{complexity unit}
\newacronym{mse}{MSE}{mean squared error}
\newacronym{relu}{ReLU}{rectified linear unit}
\newacronym{scg}{SCG}{scaled conjugate gradient}
\newacronym{sgd}{SGD}{stochastic gradient descent}
\makenoidxglossaries

\newcommand{\vect}[1]{\mathbf{#1}}

\begin{document}

\title{A Comparison of Stochastic Gradient Descent, Scaled Conjugate Gradient and LeapFrog for the Training of Feedforward Neural Networks}

\author{\IEEEauthorblockN{S.L. Johnson}
\IEEEauthorblockA{\textit{27356485}\\
\textit{Department of Computer Science}\\
\textit{Stellenbosch University}\\
27356485@sun.ac.za\\
Code \url{https://github.com/FearlessSnail/ML-assignment3.git}}
}

\maketitle

\begin{abstract}
This study compared three algorithms for the training of feedforward neural
networks with one hidden layer: stochastic gradient descent, scaled conjugate
gradient and the LeapFrog dynamic method. Five classification and four function
approximation problems of varying complexity served as benchmarks. After the
optimal number of hidden units and the best control parameters had been
determined empirically for each problem, the algorithms were compared over 30
paired runs per problem under an equal budget of computational work. No algorithm generalised significantly better across the
nine problems. Stochastic gradient descent attained a validation cost common to all three
algorithms with significantly less computational work, whereas the two batch methods required
significantly less wall-clock time. Scaled conjugate gradient was insensitive to
the control parameters over the grid tested and was concluded to be the most practical of the
three algorithms.
\end{abstract}

\section{Introduction}

A feedforward neural network learns a mapping from inputs to outputs through the
adjustment of the weights of the network. The adjustment is an unconstrained
optimisation problem, and the choice of optimisation algorithm affects both the
quality of the trained network and the cost of training. The present study
compares three algorithms that rest on different principles. \Gls{sgd} follows
the negative gradient of the error on small random subsets of the data, whereas the
\gls{scg} algorithm of M{\o}ller~\cite{moller1993} uses second-order information
to choose the search direction and the step size. The LeapFrog algorithm of
Snyman~\cite{snyman1982,snyman1983} models the weights as a particle that moves
under the force of the negative gradient. \Gls{scg} and LeapFrog are batch methods,
because every step uses the gradient over the full training set.

A fair comparison requires care: each algorithm has to receive the same
architecture, the best control parameters for the problem and an equal amount of
computational work. The control parameters of an algorithm are the constants that
the user sets before training, such as the learning rate. Because every training
run also depends on a random data split and on random initial weights, the results
have to be tested statistically. The purpose of the study was to identify the best
of the three algorithms under such a controlled comparison. After the optimal number of hidden
units and the best control parameters had been determined for nine benchmark
problems, the algorithms were compared on 30 fresh paired runs per problem with
rank-based statistical tests.

No algorithm generalised significantly better across the nine problems. \Gls{sgd}
learned significantly faster per unit of computational work, whereas the two batch
methods were significantly faster in wall-clock time. \Gls{scg} was also
insensitive to the control parameters over the grid tested.

The remainder of the report is structured as follows. Section~II presents the
background on neural networks, the three algorithms and the statistical tests,
and Section~III describes the implementation. Section~IV details the empirical
process, Section~V presents and discusses the results, and Section~VI concludes
the study.

\section{Background}

This section provides the theoretical foundation for the study. Feedforward
neural networks and generalisation are described first, followed by the three
training algorithms and a common measure of computational work. The
section closes with the statistical tests used to compare algorithms.

\subsection{Feedforward Neural Networks}

A feedforward neural network with one hidden layer maps an input vector
$\vect{x}$ of $d$ components to an output vector $\vect{y}$ of $c$ components
through $h$ hidden units~\cite{engelbrecht2007}. Each unit passes a weighted sum
of the values of the preceding layer, referred to as the net input, through an
activation function. A bias is the weight of an additional input fixed at one.
Let $\vect{U}$ denote the $(d+1) \times h$ matrix of weights from the inputs to
the hidden units, and $\vect{V}$ the $(h+1) \times c$ matrix of weights from the
hidden units to the outputs. The network computes
\begin{equation}
\vect{z} = f\big(\vect{U}^{\top} [\vect{x}; 1]\big), \qquad
\vect{y} = g\big(\vect{V}^{\top} [\vect{z}; 1]\big),
\end{equation}
where $\vect{z}$ denotes the vector of hidden unit outputs and $[\vect{x}; 1]$
denotes $\vect{x}$ extended with a constant one. The functions $f$ and $g$ denote
the hidden and output activation functions, applied to each component.

The classical activation function is the logistic sigmoid
$f(\nu) = 1/(1 + \exp(-\nu))$, where $\nu$ denotes the net input of a unit. The
sigmoid is smooth and bounded on the interval $(0, 1)$, whereas the \gls{relu},
given by $f(\nu) = \max(0, \nu)$, is a piecewise linear alternative. A linear
output function suits function approximation with unbounded targets, whereas a
sigmoid output function suits classification by keeping each output in the
interval $(0, 1)$. Hornik \emph{et al.}~\cite{hornik1989} showed
that one hidden layer with sufficient sigmoid units approximates any continuous
function on a bounded domain to any desired accuracy. The number of hidden units
therefore determines the capacity of the network: too few units are unable to
represent the function, and too many fit the noise in the training
data~\cite{engelbrecht2007}.

All $q = (d+1)h + (h+1)c$ weights form a single weight vector $\vect{w}$. Training
minimises the squared error, summed over the outputs and averaged over the $n$
training patterns, given by
\begin{equation}
e(\vect{w}) = \frac{1}{2n} \sum_{i=1}^{n} \lVert \vect{y}_i - \vect{t}_i \rVert^2,
\end{equation}
where $\vect{y}_i$ and $\vect{t}_i$ denote the output and the target for pattern
$i$. The cost as a function of the weights is referred to as the error surface. A
forward pass computes the outputs of all units for every training pattern, and
backpropagation adds a backward pass that propagates the output errors back
through the weights to give the gradient
$\nabla e(\vect{w})$~\cite{rumelhart1986}.

\subsection{Generalisation and Overfitting}

The purpose of training is a low error on unseen patterns, referred to as the
generalisation error~\cite{engelbrecht2007}. The training error decreases as
training continues and as units are added, whereas the generalisation error
eventually increases again because the network starts to fit noise. The data are
therefore divided into three disjoint sets. The cost is minimised on the training
set and the validation set guides every choice. The test set measures the final
network once, so that the choices do not bias the estimate of the generalisation
error. The training cost and
the validation cost are the cost function $e$ evaluated on the training set and on
the validation set. Early stopping limits overfitting by retaining the weights with the
lowest validation cost seen during training~\cite{prechelt1998}.

The validation error is the generalisation error measured on the validation
set. A candidate configuration, such as a number of hidden units, is judged on the
mean validation error over repeated runs. The standard error of the mean is the sample
standard deviation divided by the square root of the number of runs. Because
the candidate with the lowest mean tends to be a lucky and needlessly complex
winner, the one-standard-error rule selects the simplest candidate with a mean
within one standard error of the lowest mean~\cite{hastie2009}. The standard
error in the rule is that of the candidate with the lowest mean.

\subsection{Stochastic Gradient Descent}

\Gls{sgd} moves the weights in the direction of the negative gradient, with the
full gradient replaced by the gradient over a small random batch of patterns. A pass over the training set is referred to as an
epoch, and batches of $b$ patterns give $n/b$ updates per epoch. A momentum term
accumulates the previous updates~\cite{polyak1964,rumelhart1986}. Let $B$ denote
the current batch, $e_B$ the cost over $B$ and $k$ the iteration index. The update
is given by
\begin{equation}
\Delta\vect{w}_k = \gamma \, \Delta\vect{w}_{k-1} - \eta \, \nabla e_B(\vect{w}_k),
\qquad \vect{w}_{k+1} = \vect{w}_k + \Delta\vect{w}_k,
\end{equation}
where $\Delta\vect{w}_k$ denotes the weight change, $\eta$ the learning rate and
$\gamma$ the momentum. A small learning rate slows progress, and a large learning
rate causes divergence. The batch size trades the number of updates against how closely each batch
gradient matches the full gradient.
Because \gls{sgd} does not adapt the step size, the performance depends directly
on $\eta$ and $\gamma$.

\subsection{Scaled Conjugate Gradient}

The Hessian matrix $\vect{H}$ contains the second derivatives of $e$ with respect
to the weights. Conjugate gradient methods minimise a quadratic function of $q$
variables in at most $q$ steps~\cite{hestenes1952}. Step $k$ searches along a
direction $\vect{p}_k$ conjugate to every earlier direction $\vect{p}_j$, that is
$\vect{p}_k^{\top} \vect{H} \vect{p}_j = 0$ for $j < k$, so that a minimisation
along a conjugate direction does not undo the minimisation along earlier
directions. The error surface of a neural network is not quadratic, however. The
Hessian is positive definite when the error curves upward along every direction, a
property often lost far from a minimum. A step size computed from the curvature is then
unreliable, so classical conjugate gradient methods determine the step size
through a line search, a one-dimensional minimisation of
the error along the search direction. M{\o}ller~\cite{moller1993}
noted that the line search costs additional evaluations of the error or the gradient in every iteration.

M{\o}ller~\cite{moller1993} proposed \gls{scg} to remove the line search. First,
the product of the Hessian with the search direction is approximated through a
finite difference of two gradients, given by
\begin{equation}
\vect{s}_k = \frac{\nabla e(\vect{w}_k + \sigma_k \vect{p}_k) - \nabla e(\vect{w}_k)}{\sigma_k},
\end{equation}
where $\sigma_k = \sigma / \lVert \vect{p}_k \rVert$ for a small constant $\sigma$.
The curvature along the search direction is $\vect{p}_k^{\top} \vect{H}
\vect{p}_k$, approximated by $\vect{p}_k^{\top} \vect{s}_k$. Second, a scale
parameter $\lambda_k$ is added to the curvature. As in the Levenberg--Marquardt
method, the addition acts as if a multiple of the identity matrix were added to
the Hessian, and a larger $\lambda_k$ gives a shorter step. The scaled curvature is
\begin{equation}
\kappa_k = \vect{p}_k^{\top} \vect{s}_k + \lambda_k \lVert \vect{p}_k \rVert^2.
\end{equation}
The scale parameter is raised whenever $\kappa_k \le 0$, so that the quadratic
model always has a minimum along the search direction. Third, a comparison parameter $\rho_k$ measures
the agreement between the actual reduction of the error and the reduction predicted
by a quadratic model, the second-order Taylor approximation of the error along the
search direction. A large $\rho_k$ indicates a reliable model and lowers
$\lambda_k$ to allow longer steps, whereas a small $\rho_k$ raises $\lambda_k$ and
shortens the steps.

Algorithm~\ref{alg:scg} summarises \gls{scg}, with the negative gradient
$\vect{r}_k = -\nabla e(\vect{w}_k)$, the step size $\alpha_k = \mu_k / \kappa_k$
and $\mu_k = \vect{p}_k^{\top} \vect{r}_k$. A temporary scale parameter
$\bar\lambda_k$ records the scaling already added to the curvature, so that a
rejected step is retried without computing $\vect{s}_k$ again. Each iteration costs two
gradient evaluations and one error evaluation. M{\o}ller~\cite{moller1993}
recommended $\sigma \le 10^{-4}$ and an initial scale parameter
$\lambda_1 \le 10^{-6}$. Having found the performance on the 5-bit parity
problem not significantly affected by $\sigma \le 10^{-4}$, M{\o}ller concluded
that \gls{scg} has no user-dependent parameters crucial to success.

\begin{algorithm}[t]
\caption{Scaled conjugate gradient}
\label{alg:scg}
\begin{algorithmic}[1]
\STATE Set $\vect{p}_1 = \vect{r}_1$, $\bar\lambda_1 = 0$, $k = 1$ and mark the last step as successful
\REPEAT
\IF{the last step was successful}
\STATE Compute $\vect{s}_k$ and the curvature $\vect{p}_k^{\top} \vect{s}_k$
\ENDIF
\STATE Add $(\lambda_k - \bar\lambda_k) \lVert \vect{p}_k \rVert^2$ to the curvature to obtain $\kappa_k$
\STATE Raise $\lambda_k$ until the curvature is positive, if $\kappa_k \le 0$
\STATE Set $\alpha_k = \mu_k / \kappa_k$ and compute $\rho_k = 2 \kappa_k [e(\vect{w}_k) - e(\vect{w}_k + \alpha_k \vect{p}_k)] / \mu_k^2$
\IF{$\rho_k \ge 0$}
\STATE Accept $\vect{w}_{k+1} = \vect{w}_k + \alpha_k \vect{p}_k$ and set $\bar\lambda_k = 0$
\STATE Set $\beta_k = (\lVert \vect{r}_{k+1} \rVert^2 - \vect{r}_{k+1}^{\top} \vect{r}_k) / \mu_k$ and $\vect{p}_{k+1} = \vect{r}_{k+1} + \beta_k \vect{p}_k$, or $\vect{p}_{k+1} = \vect{r}_{k+1}$ after every $q$ iterations
\STATE Divide $\lambda_k$ by four if $\rho_k \ge 0.75$
\ELSE
\STATE Reject the step, set $\bar\lambda_k = \lambda_k$ and mark the step as unsuccessful
\ENDIF
\STATE Increase $\lambda_k$ if $\rho_k < 0.25$, and increase $k$ by one
\UNTIL{the gradient vanishes or the budget is spent}
\end{algorithmic}
\end{algorithm}

\subsection{LeapFrog}

Snyman~\cite{snyman1982} treated minimisation as a problem of mechanics. The
weights are the position of a particle of unit mass, the error $e(\vect{w})$ is the
potential energy of the particle, and the force on the particle is the negative
gradient. The particle therefore follows $\ddot{\vect{w}} = -\nabla e(\vect{w})$,
where the double dot denotes the second derivative with respect to time. The
kinetic energy of the particle is $\frac{1}{2} \lVert \dot{\vect{w}} \rVert^2$,
where the single dot denotes the first derivative. Because the force
field is conservative, the sum of the potential and the kinetic energy is constant,
and a decrease in the error is matched by an increase in the kinetic energy. The trajectory is integrated with the leap-frog scheme, given by
\begin{equation}
\vect{w}_{k+1} = \vect{w}_k + \tau \vect{v}_k, \qquad
\vect{v}_{k+1} = \vect{v}_k + \tau \vect{a}_{k+1},
\end{equation}
where $\vect{v}_k$ denotes the velocity, $\vect{a}_k = -\nabla e(\vect{w}_k)$ the
acceleration and $\tau$ the time step.

An undisturbed particle would oscillate indefinitely, so LeapFrog interferes
whenever the particle loses speed, a sign that the particle climbs. The particle then restarts from the midpoint of the last
step with a reduced velocity, or from rest after repeated restarts. The original
algorithm, LFOP1, limits the step length $\tau \lVert \vect{v}_k \rVert$ to a
maximum $\delta$~\cite{snyman1982}, and the improved version LFOP1(b) adds an
automatic control of the time step~\cite{snyman1983}. The force reverses when the new acceleration
points against the previous one, that is when
$\vect{a}_{k+1}^{\top} \vect{a}_k \le 0$. After every step shorter than
$\delta$ and without a reversal, the time step is multiplied by a growth factor that
increases by $\delta_1$ after each such step and returns to one after a reversal.
The time step halves after $m$ consecutive reversals. Each step of LFOP1(b),
summarised in Algorithm~\ref{alg:leapfrog}, costs one gradient evaluation. Snyman used $\tau_0 = 0.5$ and $\delta = 1$ for
LFOP1~\cite{snyman1982}, and added $m = 3$ and $\delta_1 = 0.001$ for
LFOP1(b)~\cite{snyman1983}.

\begin{algorithm}[t]
\caption{LeapFrog LFOP1(b)}
\label{alg:leapfrog}
\begin{algorithmic}[1]
\STATE Set $\vect{v}_0 = \frac{1}{2} \tau_0 \vect{a}_0$ and $\tau = \tau_0$
\REPEAT
\IF{the step length $\tau \lVert \vect{v}_k \rVert$ exceeds $\delta$}
\STATE Scale $\vect{v}_k$ so that the step length equals $\delta$
\ELSE
\STATE Increase the growth factor by $\delta_1$ and multiply $\tau$ by the factor
\ENDIF
\IF{the force reversed direction in $m$ consecutive steps}
\STATE Halve $\tau$, move to the midpoint of the last step and damp the velocity
\ENDIF
\STATE Set $\vect{w}_{k+1} = \vect{w}_k + \tau \vect{v}_k$ and $\vect{v}_{k+1} = \vect{v}_k + \tau \vect{a}_{k+1}$
\IF{$\lVert \vect{v}_{k+1} \rVert \le \lVert \vect{v}_k \rVert$}
\STATE Restart from the midpoint with the velocity $\frac{1}{4}(\vect{v}_{k+1} + \vect{v}_k)$, or with zero velocity after repeated restarts
\ENDIF
\STATE Increase $k$ by one
\UNTIL{the gradient vanishes or the budget is spent}
\end{algorithmic}
\end{algorithm}

\subsection{Computational Work}

Iterations of the three algorithms are not comparable, so M{\o}ller~\cite{moller1993}
measured work in \glspl{cu}. One \gls{cu} is one forward pass over the
full training set, and a gradient over the full training set costs three
\glspl{cu}. One \gls{sgd} epoch therefore costs three \glspl{cu} regardless of the
batch size, as does one LeapFrog step, while one \gls{scg} iteration costs about
seven.

\subsection{Statistical Comparison of Algorithms}

The data split, the initial weights and the order of the batches are all random,
so algorithms are compared over repeated runs. In a paired design, one seed fixes
the data split and the initial weights for every algorithm, so that the variation
between splits cancels in the differences. The paired $t$-test assumes
normally distributed differences, an assumption that the Shapiro--Wilk test
assesses~\cite{shapiro1965}. The Wilcoxon signed-rank test~\cite{wilcoxon1945}
makes no such assumption~\cite{demsar2006} and compares the rank sums of the
positive and the negative differences. The Hodges--Lehmann estimator, the median of all pairwise
averages of the differences, estimates the size and direction of the difference between two
algorithms~\cite{hodges1963}.
The procedure of Holm~\cite{holm1979} adjusts the probability values of multiple
tests to control the probability of any false positive.

Dem\v{s}ar~\cite{demsar2006} recommended rank-based tests for the comparison of
algorithms over multiple problems. The algorithms are ranked on each problem, and the Friedman
test~\cite{friedman1937} assesses whether the mean ranks differ.
The Nemenyi test~\cite{demsar2006,nemenyi1963} declares two algorithms different when the mean
ranks differ by more than a critical difference. For three algorithms and nine
problems, the critical difference at a significance level of 0.05 is 1.10 mean
ranks~\cite{demsar2006}.

In summary, \gls{sgd} relies on the gradient alone, \gls{scg} adds curvature
information without a line search, and LeapFrog follows a damped physical
trajectory. \Glspl{cu} make the work of the three algorithms comparable, and
rank-based tests compare the algorithms without an assumption of normality.

\section{Implementation}

This section describes the implementation, starting with the software and
followed by the network and the training algorithms.

\subsection{Software}

The implementation is written in Python\footnote{\url{https://www.python.org}},
and the NumPy library\footnote{\url{https://numpy.org}} provides all numerical
computation,
including the network and the three training algorithms. The training algorithms
are implemented from the source papers rather than taken from a library, so that
every algorithm counts work in exactly the same way and is verifiable against the
source paper. The scikit-learn
library\footnote{\url{https://scikit-learn.org}} supplies five benchmark problems
and the stratified division of the data but plays no part in training. The SciPy
library\footnote{\url{https://scipy.org}} provides the statistical tests, and the
Matplotlib library\footnote{\url{https://matplotlib.org}} produces the figures.

\subsection{Network and Training Algorithms}

The network stores all weights in one vector, so the training algorithms treat
the network as a function from the weights to the cost and the gradient. A
comparison with central finite differences confirmed the gradient to a relative
error of $4.4 \times 10^{-8}$. The initial weights are drawn uniformly from the
interval $[-1/\sqrt{\phi}, 1/\sqrt{\phi}]$, where $\phi$ denotes the number of
inputs to a unit, bias included~\cite{engelbrecht2007}. The interval keeps the initial net inputs in the
responsive range of the sigmoid~\cite{lecun1998}. Every evaluation of the cost or
the gradient adds the \glspl{cu} of the evaluation to one counter, so that the
work of all three algorithms is counted in the same way.

The three algorithms share one interface: each receives the initial
weights, the cost function and a monitor. \Gls{sgd} reshuffles the patterns every
epoch, while the implementations of \gls{scg} and LeapFrog follow the steps of the source
papers one by one. \Gls{scg} additionally stops when the search direction vanishes
or the scale parameter exceeds $10^{20}$, two cases that the paper does not
discuss. The monitor evaluates the training and the validation cost at the end of an
iteration, at most once every three \glspl{cu}. The monitor records the weights with the
lowest validation cost and ends the
run when the budget is spent. One seed fixes the data split, the initial weights and the order of the \gls{sgd}
batches, so that runs with the same seed are paired and every result is
reproducible.

\section{Empirical Process}

This section describes the empirical process. The benchmark problems and the
preprocessing of the data come first, followed by the selection of the number of
hidden units and of the control parameters. The section closes with the
performance criteria and the comparison protocol.

\subsection{Benchmark Problems}

Table~\ref{tab:problems} lists the nine benchmark problems, chosen to vary in
size, dimension, number of classes, non-linearity and noise.

\begin{table}[t]
\centering
\caption{Benchmark problems}
\label{tab:problems}
\footnotesize
\begin{tabular}{@{}llrrrr@{}}
\toprule
Problem & Task & $n$ & $d$ & $c$ & Budget \\
\midrule
Iris & classification & 150 & 4 & 3 & 600 \\
Wine & classification & 178 & 13 & 3 & 900 \\
Cancer & classification & 569 & 30 & 2 & 900 \\
Digits & classification & 1797 & 64 & 10 & 2400 \\
Spirals & classification & 400 & 2 & 2 & 6000 \\
Sinc & regression & 300 & 1 & 1 & 3000 \\
Sin2D & regression & 600 & 2 & 1 & 4000 \\
Friedman1 & regression & 600 & 10 & 1 & 2000 \\
Diabetes & regression & 442 & 10 & 1 & 900 \\
\bottomrule
\end{tabular}

\end{table}

Five problems are classification problems. Iris contains four measurements of 150
flowers of three species. Wine contains 13 chemical measurements
of 178 wines from three
cultivars\footnote{\url{https://archive.ics.uci.edu/dataset/109/wine}}, on scales
that differ by four orders of magnitude. Cancer contains 30 features of cell nuclei
from 569 breast tissue samples, labelled malignant or
benign\footnote{\url{https://archive.ics.uci.edu/dataset/17}}. Digits contains
1\,797 images of handwritten digits of eight by eight pixels, in 10 classes, in the
version distributed with scikit-learn, the test portion of the original data
set\footnote{\url{https://archive.ics.uci.edu/dataset/80}}. Spirals is a version of the
two-spirals problem~\cite{lang1988}, with 400 points on two interlocking spirals of
1.25 turns separated by a highly non-linear boundary.

Four problems are function approximation problems. Let $x_1$ and $x_2$ denote the
first two components of the input vector. Sinc samples $\sin(x_1)/x_1$ at 300
points on the interval $[-10, 10]$. Sin2D samples $\sin(2\pi x_1)\cos(2\pi x_2)$
at 600 points on the unit square. Both add Gaussian noise with a standard deviation
of 0.05. Friedman1 samples the 10-input test function of Friedman~\cite{friedman1991} at
600 points, with five irrelevant inputs and Gaussian noise of standard deviation
one. Diabetes contains 10 measurements of 442 patients and the
progression of the disease after one year.

\subsection{Data Preprocessing and Network Configuration}

Every run divided the data into 60 percent training, 20 percent validation and 20
percent test patterns, stratified for classification so that each set kept the
class proportions. Every input was standardised to zero mean and
unit variance with the statistics of the training set alone. Inputs constant on
the training set, such as pixels of Digits that are blank in every training image, were centred
but not scaled.
Standardisation prevents large inputs from saturating the sigmoid
units~\cite{engelbrecht2007} and improves the conditioning of the error
surface~\cite{lecun1998}. The conditioning is the ratio of the largest to the
smallest curvature of the error surface, given by the eigenvalues of the
Hessian~\cite{lecun1998}. A
poorly conditioned surface forms long narrow valleys, where no single step size
suits every direction. The effect of a weight
change on the net input depends on the scale of the inputs. Standardisation
therefore also gives the learning rate of \gls{sgd} and the step limit of LeapFrog
a comparable effect across problems.

Classification targets were coded with one output per class, at 0.9 for the true
class and 0.1 for the others. Because a sigmoid attains zero or one only in the limit, targets of zero and
one would drive the weights towards infinity~\cite{engelbrecht2007,lecun1998}. No problem contained missing
values, and the only nominal input was the binary sex variable of Diabetes.

Every network had one hidden layer of sigmoid units, because \gls{scg} and LeapFrog
both assume a smooth cost function. The output units were sigmoid for
classification and linear for function approximation, and every problem used the
averaged sum of squared errors as the cost function. Every run on a problem received
the same budget of \glspl{cu}, listed in Table~\ref{tab:problems}. On seven
problems, the budgets left ample room beyond the typical point of convergence: over
the final tenth of the budget, the median improvement of the lowest validation cost
over the comparison runs was zero. On Sinc and Sin2D, the median improvement over
the final tenth was up to 3 percent, so the results on the two problems reflect the
budget as well as the algorithms.

\subsection{Selection of the Number of Hidden Units}

For each size, the validation error was averaged over repeated runs. The size with
the lowest mean set a threshold one standard error above that mean, and the optimal
number of hidden units $h^*$ was the smallest size with a mean at or below the
threshold. For
classification, the validation error is the fraction of misclassified validation
patterns. For
function approximation, the validation error is the \gls{mse} on the validation set.
Early stopping and the tie-break of Section~IV-D use the validation cost instead.
The validation error measures generalisation, and the mean over repeated runs
estimates the expected error. The smallest size at or below the threshold avoids a choice
based on differences below the level of noise.

Stage A trained networks of 1, 2, 4, 6, 8, 12, 16, 24 and 32 hidden units, a grid
dense at the small sizes where the error changes fastest. Each algorithm trained
every size with five seeds, at the default control parameters of
Table~\ref{tab:grids}, because the control parameters were unknown before the
architecture was chosen. The validation errors of the three algorithms were
pooled, so that one size served all three algorithms and the comparison held the
architecture fixed. Stage A2 repeated stage A with the
control parameters fixed at the values that stage B tuned at $h^*$, to test
whether the choice depended on the order of the
two stages.

\subsection{Selection of the Control Parameters}

The best control parameters of an algorithm on a problem were chosen with the
same rule. At $h^*$, the setting with the lowest mean validation error set a
threshold one standard error above that mean. Among the settings at or below the
threshold, the setting with the least mean work to the lowest validation cost of a
run was preferred. Stage B evaluated every combination in
Table~\ref{tab:grids} with 10 seeds on every problem. A full grid was necessary because the
parameters interact: the effective step of \gls{sgd}, for example, grows
approximately as $\eta/(1-\gamma)$. Each range extended from conservative to
aggressive values, on an approximately logarithmic scale for $\eta$, $b$, $\sigma$,
$\lambda_1$, $\tau_0$, $\delta$ and $\delta_1$. A first pass with five seeds selected a setting of \gls{sgd} close to divergence on Sin2D,
so the number of seeds was raised to 10 for every problem and algorithm.

\begin{table}[t]
\centering
\caption{Grids of control parameters}
\label{tab:grids}
\footnotesize
\begin{tabular}{@{}llll@{}}
\toprule
Algorithm & Parameter & Values & Default \\
\midrule
\multirow{3}{*}{SGD} & $\eta$ & 0.05, 0.1, 0.5, 1, 2 & 0.1 \\
 & $\gamma$ & 0, 0.5, 0.9 & 0.9 \\
 & $b$ & 8, 32, 128 & 32 \\
\midrule
\multirow{2}{*}{SCG} & $\sigma$ & $10^{-6}$, $10^{-5}$, $10^{-4}$, $10^{-3}$ & $10^{-4}$ \\
 & $\lambda_1$ & $10^{-8}$, $10^{-6}$, $10^{-4}$ & $10^{-6}$ \\
\midrule
\multirow{4}{*}{LeapFrog} & $\tau_0$ & 0.1, 0.5, 1 & 0.5 \\
 & $\delta$ & 0.5, 2, 10, 50 & 1 \\
 & $m$ & 3, 5, 8 & 3 \\
 & $\delta_1$ & 0.001, 0.01 & 0.001 \\
\bottomrule
\end{tabular}
\end{table}

\subsection{Performance Criteria and Comparison Protocol}

Four performance criteria were defined before the comparison. The primary
criterion was the generalisation error on the test set. For classification, the
generalisation error is the misclassification rate, equal to one minus the
accuracy, where the accuracy is the fraction of test patterns classified correctly.
The \gls{mse} alone is a poor measure of accuracy for
classification~\cite{engelbrecht2007}. For
function approximation, the generalisation error is the \gls{mse} on the test set.
The second criterion was the \gls{mse} on the test set for classification as well,
a finer measure than the misclassification rate. The third criterion was the work
to target, the number of \glspl{cu} a run required to attain the target of the
split. On each split, every algorithm attained a lowest validation cost, and the target
of the split was the highest of the three. A target per split adapts the
criterion to the difficulty of the split, and every run attains the target by
construction. The fourth criterion was the wall-clock time for the full budget.

Stage C compared the algorithms on 30 seeds that no earlier stage had used, with
each algorithm at $h^*$, the tuned control parameters and the budget of the problem.
Stages A, B, A2 and C comprised 1\,215, 11\,610, 1\,215 and 810 runs. Within each
problem, the Wilcoxon signed-rank test compared every pair of algorithms, with the
Holm correction over the three pairs and a significance level of 0.05. The
rank-based test was preferred to the paired $t$-test because the rank-based test
does not assume normality~\cite{demsar2006}. The Shapiro--Wilk test supported the
choice, rejecting normality of the paired differences in 16 of the 27 comparisons
of generalisation error. Differences in misclassification rate are discrete and
often tied, a property that alone departs from normality. Across the nine problems,
the Friedman test and the Nemenyi critical difference compared the mean ranks,
with each algorithm ranked on the mean of each criterion over the 30 runs. The
complete process was also repeated with \glspl{relu} in the hidden layer, with the
same grids of control parameters and a grid of sizes extended by 48, 64, 96 and
128 units. The repetition tested whether the conclusions survived a cost function
that is not smooth.

In summary, the number of hidden units and the control parameters were chosen on
validation data with the one-standard-error rule and full grids. The comparison
used 30 fresh paired runs per problem, four criteria and rank-based tests.

\section{Results and Discussion}

This section presents and discusses the results, beginning with the optimal
numbers of hidden units and the tuned control parameters. The generalisation
error, the speed of learning and the overall comparison follow.

\subsection{Optimal Number of Hidden Units}

Figure~\ref{fig:hidden} shows the mean validation error against the number of
hidden units, and Table~\ref{tab:hidden} lists the chosen sizes. The error fell
steeply while the network was too small and then flattened. On Iris, Spirals, Sinc
and Friedman1, the error rose again for the largest networks.

\begin{figure*}[!t]
\centering
\includegraphics[width=0.75\textwidth]{legend_hidden}\\[-5pt]
\subfloat[Iris]{\includegraphics{hidden_iris}}\hfil
\subfloat[Wine]{\includegraphics{hidden_wine}}\hfil
\subfloat[Cancer]{\includegraphics{hidden_cancer}}\\[-5pt]
\subfloat[Digits]{\includegraphics{hidden_digits}}\hfil
\subfloat[Spirals]{\includegraphics{hidden_spirals}}\hfil
\subfloat[Sinc]{\includegraphics{hidden_sinc}}\\[-5pt]
\subfloat[Sin2D]{\includegraphics{hidden_sin2d}}\hfil
\subfloat[Friedman1]{\includegraphics{hidden_friedman}}\hfil
\subfloat[Diabetes]{\includegraphics{hidden_diabetes}}
\caption{Mean validation error against the number of hidden units}
\label{fig:hidden}
\end{figure*}

\begin{table}[t]
\centering
\caption{Optimal number of hidden units}
\label{tab:hidden}
\footnotesize
\begin{tabular}{@{}lrrrrrr@{}}
\toprule
& \multicolumn{5}{c}{Stage A} & A2 \\ \cmidrule(lr){2-6}\cmidrule(l){7-7}Problem & $h^*$ & Pooled & SGD & SCG & LeapFrog & $h^*$ \\
\midrule
Iris & \textbf{2} & 4 & 4 & 2 & 2 & 2 \\
Wine & \textbf{6} & 6 & 8 & 6 & 6 & 6 \\
Cancer & \textbf{1} & 1 & 1 & 16 & 12 & 16 \\
Digits & \textbf{24} & 32 & 32 & 32 & 24 & 24 \\
Spirals & \textbf{16} & 16 & 16 & 32 & 32 & 16 \\
Sinc & \textbf{6} & 12 & 12 & 8 & 12 & 8 \\
Sin2D & \textbf{12} & 12 & 16 & 12 & 12 & 12 \\
Friedman1 & \textbf{8} & 8 & 6 & 8 & 8 & 8 \\
Diabetes & \textbf{8} & 32 & 16 & 8 & 24 & 8 \\
\bottomrule
\end{tabular}

\end{table}

The chosen sizes reflected the complexity of the problems. The mean validation
error of Cancer varied by less than one misclassified validation pattern across all
sizes, so one unit sufficed for
the nearly linearly separable problem. Iris required two units, with an error of
0.136 for one unit against 0.018 for two. Every output of a network
with one hidden unit is a monotone function of the same net input. The middle
species of Iris, however, requires a high output between two low outputs, a
pattern that no monotone function produces. Spirals required 16 units and Digits 24. The individual
algorithms preferred different sizes on the flat curves: on Diabetes, the lowest
mean occurred at 16 units for \gls{sgd}, eight for \gls{scg} and 24 for LeapFrog.
Pooling the algorithms reduced the influence of noise on the choice.

Stage A2 confirmed the choice on seven of the nine problems, and both changes
occurred on flat curves. On Sinc, the choice moved from six to eight units because
six units missed the threshold by 0.001. On Cancer, the choice moved from one to 16
units, with every size performing equally well. The comparison kept the sizes of
stage A. Stage A2 also exposed a weakness of \gls{sgd}: on Sin2D, Friedman1 and
Diabetes, the learning rate tuned at $h^*$ failed on larger networks
(Figure~\ref{fig:stability}(b)).

Figure~\ref{fig:stability} explains the failure. Let $\omega$ denote the largest
eigenvalue of the Hessian with respect to the output weights. On a quadratic
function, gradient descent with momentum converges only if
$\eta\,\omega < 2(1+\gamma)$~\cite{polyak1964}, so the stability ratio
$\eta\,\omega / [2(1+\gamma)]$ must stay below one. For a linear output, the
Hessian with respect to the output weights is the mean outer product of the hidden
outputs, so $\omega$ grows as hidden units are added.
Figure~\ref{fig:stability}(a) shows the stability ratio at the initial weights
for the tuned parameters of \gls{sgd}, and Figure~\ref{fig:stability}(b) shows the
validation error of tuned \gls{sgd} in stage A2, with \gls{scg} dotted. On Sin2D and
Friedman1, the ratio exceeded one only at 32 units, where the error of \gls{sgd}
rose to the level of an untrained network. On Diabetes, the ratio already attained
one at eight units, and the error of \gls{sgd} rose from 12 units onwards.
\Gls{scg} and LeapFrog adapt the step to the error surface and therefore did not
suffer from the change of size.

\begin{figure}[t]
\centering
\subfloat[Stability ratio]{\includegraphics{stability_ratio}}\hfil
\subfloat[Validation error]{\includegraphics{stability_error}}
\caption{Stability of tuned stochastic gradient descent}
\label{fig:stability}
\end{figure}

\subsection{Control Parameters}

Table~\ref{tab:params} lists the chosen control parameters, and the column
labelled tied counts the settings at or below the threshold of Section~IV-D. A high count
indicates insensitivity to the parameters. Table~\ref{tab:sensitivity} measures sensitivity
in two ways, each as the median and the maximum over the nine problems. The first
fixes one parameter at a value and tunes the others. The parameter-fixed columns give how far the
best mean validation error attainable at the worst such value lies above the
overall best. The second is the spread between the best and the worst setting of
the whole grid. A large spread with small values for single parameters indicates
an interaction between parameters. For classification, the values are differences
in misclassification rate. For function approximation, the values are differences
in \gls{mse} on standardised targets, so a spread above one means that the worst
setting did worse than predicting the mean.

\begin{table*}[!t]
\centering
\caption{Chosen control parameters}
\label{tab:params}
\footnotesize
\begin{tabular}{@{}l rrrr rrr rrrrr@{}}
\toprule
& \multicolumn{4}{c}{SGD} & \multicolumn{3}{c}{SCG} & \multicolumn{5}{c}{LeapFrog} \\\cmidrule(lr){2-5}\cmidrule(lr){6-8}\cmidrule(l){9-13}Problem & $\eta$ & $\gamma$ & $b$ & Tied & $\sigma$ & $\lambda_1$ & Tied & $\tau_0$ & $\delta$ & $m$ & $\delta_1$ & Tied \\
\midrule
Iris & 2 & 0.9 & 8 & 4/45 & $10^{-3}$ & $10^{-4}$ & 12/12 & 1 & 50 & 8 & 0.01 & 14/72 \\
Wine & 1 & 0.9 & 8 & 33/45 & $10^{-3}$ & $10^{-6}$ & 12/12 & 0.1 & 10 & 8 & 0.01 & 14/72 \\
Cancer & 2 & 0.5 & 8 & 14/45 & $10^{-3}$ & $10^{-6}$ & 12/12 & 1 & 2 & 8 & 0.01 & 6/72 \\
Digits & 2 & 0.5 & 8 & 28/45 & $10^{-5}$ & $10^{-6}$ & 12/12 & 0.5 & 10 & 5 & 0.01 & 13/72 \\
Spirals & 1 & 0.9 & 8 & 1/45 & $10^{-5}$ & $10^{-6}$ & 4/12 & 1 & 10 & 3 & 0.01 & 3/72 \\
Sinc & 0.1 & 0.9 & 8 & 3/45 & $10^{-6}$ & $10^{-8}$ & 11/12 & 0.1 & 2 & 3 & $10^{-3}$ & 17/72 \\
Sin2D & 0.5 & 0.9 & 128 & 2/45 & $10^{-3}$ & $10^{-6}$ & 6/12 & 0.1 & 2 & 8 & $10^{-3}$ & 3/72 \\
Friedman1 & 0.5 & 0.9 & 128 & 2/45 & $10^{-5}$ & $10^{-8}$ & 9/12 & 1 & 0.5 & 8 & 0.01 & 24/72 \\
Diabetes & 1 & 0.5 & 32 & 20/45 & $10^{-4}$ & $10^{-4}$ & 12/12 & 1 & 0.5 & 3 & $10^{-3}$ & 10/72 \\
\bottomrule
\end{tabular}

\end{table*}

\begin{table}[t]
\centering
\caption{Sensitivity to the control parameters}
\label{tab:sensitivity}
\footnotesize
\setlength{\tabcolsep}{4pt}
\begin{tabular}{@{}ll rr rr@{}}
\toprule
& & \multicolumn{2}{c}{Parameter fixed} & \multicolumn{2}{c}{Whole grid} \\\cmidrule(lr){3-4}\cmidrule(l){5-6}Algorithm & Parameter & Median & Maximum & Median & Maximum \\
\midrule
SGD & $\eta$ & 0.0575 & 0.9062 & 0.5800 & 1.1526 \\
 & $\gamma$ & 0.0042 & 0.0150 &  &  \\
 & $b$ & 0.0044 & 0.0233 &  &  \\
SCG & $\sigma$ & 0.0009 & 0.0033 & 0.0022 & 0.0177 \\
 & $\lambda_1$ & 0.0007 & 0.0025 &  &  \\
LeapFrog & $\tau_0$ & 0.0033 & 0.0067 & 0.0267 & 0.7463 \\
 & $\delta$ & 0.0053 & 0.0164 &  &  \\
 & $m$ & 0.0025 & 0.0147 &  &  \\
 & $\delta_1$ & 0.0033 & 0.0081 &  &  \\
\bottomrule
\end{tabular}

\end{table}

\Gls{sgd} preferred aggressive settings, with large learning rates and small
batches, and was sensitive to the settings. A
learning rate of at least 0.5 was chosen on eight problems, and a batch size of
eight, the size with the most updates per \gls{cu}, on six. On
Spirals, only one of the 45 settings lay at or below the threshold. The
learning rate was the decisive parameter, with the largest parameter-fixed value
in Table~\ref{tab:sensitivity}, 0.906 on Sin2D.
At a learning rate of two, 22.5 percent of the runs never improved on the initial
weights, so the best settings of \gls{sgd} lay close to the boundary of
divergence.

\Gls{scg} was effectively free of parameters: over the whole grid, the mean
validation error varied by at most 0.018 on any problem, and all 12 settings were
tied on five problems. The present grid confirmed the finding of
M{\o}ller~\cite{moller1993} that the performance was insensitive to $\sigma$ and
extended the finding to $\sigma = 10^{-3}$ and to
$\lambda_1$. The results of M{\o}ller already hinted at the first extension, with
$\sigma = 10^{-3}$ giving the lowest mean number of \glspl{cu} on parity. The scale
parameter adapts during training, so the initial values affect only the first
iterations.

LeapFrog was insensitive to each parameter alone, with parameter-fixed values of
at most 0.016,
but the spread over the whole grid attained 0.75 on Sinc and 0.70 on Sin2D. The
poor settings combined the largest step limit of 50 with $m = 8$, the most
reversals allowed before the time step halves. On Sinc, the combination gave a mean validation error of 0.528, against
0.028 for a step limit of 0.5. A long permitted step and slow
damping let the particle swing across the oscillations of the target function.
Snyman~\cite{snyman1983} suggested that a large step limit was advantageous on the
test functions, but on neural networks a large step limit was harmful when
combined with slow damping.

\subsection{Generalisation Error}

Table~\ref{tab:results} reports the generalisation error, and
Table~\ref{tab:pairwise-error} the pairwise tests. Each pairwise cell gives the
significantly better algorithm or a tie, and the number of runs won by each
algorithm of the pair. The standard deviations over the 30 runs were of the same size as the
differences between the means, and 20 of the 27 pairwise comparisons were ties.

\begin{table}[t]
\centering
\caption{Generalisation error over 30 runs}
\label{tab:results}
\footnotesize
\setlength{\tabcolsep}{3pt}
\begin{tabular}{@{}l rrr@{}}
\toprule
Problem & SGD & SCG & LeapFrog \\
\midrule
Iris & \textbf{0.0411 $\pm$ 0.0368} & 0.0433 $\pm$ 0.0363 & 0.0489 $\pm$ 0.0358 \\
Wine & 0.0204 $\pm$ 0.0218 & \textbf{0.0194 $\pm$ 0.0254} & 0.0269 $\pm$ 0.0225 \\
Cancer & 0.0281 $\pm$ 0.0153 & \textbf{0.0275 $\pm$ 0.0178} & 0.0289 $\pm$ 0.0148 \\
Digits & \textbf{0.0329 $\pm$ 0.0097} & 0.0344 $\pm$ 0.0092 & 0.0531 $\pm$ 0.0177 \\
Spirals & 0.0200 $\pm$ 0.0238 & \textbf{0.0117 $\pm$ 0.0164} & 0.0254 $\pm$ 0.0691 \\
\midrule
Sinc & 0.0286 $\pm$ 0.0063 & \textbf{0.0245 $\pm$ 0.0053} & 0.0251 $\pm$ 0.0054 \\
Sin2D & 0.0194 $\pm$ 0.0038 & 0.0196 $\pm$ 0.0042 & \textbf{0.0176 $\pm$ 0.0044} \\
Friedman1 & 0.0770 $\pm$ 0.0129 & 0.0977 $\pm$ 0.0663 & \textbf{0.0730 $\pm$ 0.0171} \\
Diabetes & \textbf{0.5447 $\pm$ 0.0989} & 0.5453 $\pm$ 0.0958 & 0.5475 $\pm$ 0.0897 \\
\bottomrule
\end{tabular}

\end{table}

\begin{table}[t]
\centering
\caption{Pairwise comparison of generalisation error}
\label{tab:pairwise-error}
\footnotesize
\setlength{\tabcolsep}{4pt}
\begin{tabular}{@{}l lll@{}}
\toprule
Problem & SCG, SGD & LeapFrog, SGD & LeapFrog, SCG \\
\midrule
Iris & tie 3:5 & tie 4:11 & tie 3:7 \\
Wine & tie 4:2 & tie 6:11 & tie 6:10 \\
Cancer & tie 8:7 & tie 10:11 & tie 8:13 \\
Digits & tie 9:17 & SGD 2:25 & SCG 5:25 \\
Spirals & SCG 14:2 & tie 11:11 & tie 6:13 \\
Sinc & SCG 25:5 & LeapFrog 24:6 & tie 13:17 \\
Sin2D & tie 15:15 & LeapFrog 23:7 & LeapFrog 22:8 \\
Friedman1 & tie 17:13 & tie 18:12 & tie 21:9 \\
Diabetes & tie 17:13 & tie 17:13 & tie 15:15 \\
\bottomrule
\end{tabular}

\end{table}

The seven significant differences had identifiable causes. \Gls{sgd} and \gls{scg}
beat LeapFrog on Digits, with errors of 0.033 and 0.034 against 0.053.
The median training cost of LeapFrog was 0.057 against 0.009 for \gls{sgd}, so
the cause was a failure of optimisation rather than overfitting.
Figure~\ref{fig:leapfrog} shows one run of LeapFrog on Digits and on Sin2D at the
tuned parameters. The training cost on Digits stopped improving within the first
tenth of the budget, while the time step grew to 100, against at most 1.4 on
Sin2D. With the step limit active on 45 percent of the steps, LeapFrog degenerated
into a method of fixed step length.

\begin{figure}[t]
\centering
\subfloat[Lowest training cost]{\includegraphics{leapfrog_cost}}\hfil
\subfloat[Time step]{\includegraphics{leapfrog_step}}
\caption{One LeapFrog run on Digits and Sin2D}
\label{fig:leapfrog}
\end{figure}

Both batch methods beat \gls{sgd} on Sinc, with a \gls{mse} of about 0.025 against
0.029. \Gls{sgd} used batches of eight patterns on Sinc, and ended at a median
training cost of 0.0133 against 0.0111 for \gls{scg}. The gradient of a small batch
keeps the weights moving around the minimum, so \gls{sgd} settled less precisely.
LeapFrog beat both other algorithms on Sin2D, with the lowest training cost of the
three. \Gls{scg} beat \gls{sgd} on Spirals, at 0.012 against 0.020.

\subsection{Speed of Learning and Wall-Clock Time}

Table~\ref{tab:speed} reports the median work to target and the mean wall-clock
time. \Gls{sgd} attained the target first on eight of the nine problems and won 13 of
the 18 pairwise comparisons of work that involved \gls{sgd} and tied the other
five.
Figure~\ref{fig:curves} shows the learning curves of four problems. On Digits, \gls{sgd} attained the target within two epochs, while LeapFrog
levelled off above the other two algorithms. One epoch of \gls{sgd} with batches of eight patterns
performs $n/8$ updates for three \glspl{cu}, against one update for a LeapFrog
step. Early in training, the gradient of a small batch points in nearly the same
direction as the full gradient, so each update makes almost full progress. The
batch methods started slowly and caught up later, with \gls{scg} ending lowest on
Spirals and LeapFrog on Sin2D.

\begin{table}[t]
\centering
\caption{Median work to target and mean wall-clock time}
\label{tab:speed}
\footnotesize
\setlength{\tabcolsep}{4pt}
\begin{tabular}{@{}l rrr rrr@{}}
\toprule
& \multicolumn{3}{c}{Complexity units} & \multicolumn{3}{c}{Seconds} \\\cmidrule(lr){2-4}\cmidrule(l){5-7}Problem & SGD & SCG & LeapFrog & SGD & SCG & LeapFrog \\
\midrule
Iris & \textbf{30} & 186 & 213 & 0.16 & 0.02 & \textbf{0.02} \\
Wine & \textbf{12} & 73 & 138 & 0.30 & 0.03 & \textbf{0.03} \\
Cancer & \textbf{24} & 84 & 90 & 0.94 & 0.04 & \textbf{0.04} \\
Digits & \textbf{6} & 89 & 111 & 9.38 & \textbf{2.77} & 3.09 \\
Spirals & \textbf{710} & 1383 & 1473 & 4.48 & \textbf{0.42} & 0.44 \\
Sinc & 2133 & \textbf{1836} & 2535 & 1.46 & 0.11 & \textbf{0.11} \\
Sin2D & \textbf{2499} & 3640 & 3068 & 0.49 & 0.41 & \textbf{0.35} \\
Friedman1 & \textbf{384} & 1028 & 1017 & 0.28 & 0.14 & \textbf{0.12} \\
Diabetes & \textbf{28} & 80 & 104 & 0.19 & \textbf{0.04} & 0.04 \\
\bottomrule
\end{tabular}

\end{table}

\begin{figure*}[!t]
\centering
\includegraphics[width=0.36\textwidth]{legend_curves}\\[-5pt]
\subfloat[Digits]{\includegraphics{curves_digits}}\hfil
\subfloat[Spirals]{\includegraphics{curves_spirals}}\hfil
\subfloat[Sinc]{\includegraphics{curves_sinc}}\hfil
\subfloat[Sin2D]{\includegraphics{curves_sin2d}}
\caption{Median lowest validation cost against work, with the interquartile range shaded}
\label{fig:curves}
\end{figure*}

The advantage of \gls{sgd} narrowed on Sinc and Sin2D, where every algorithm
required more than 1\,800 \glspl{cu} to attain the target. The target therefore lay
close to the lowest validation cost attainable within the budget. M{\o}ller~\cite{moller1993} noted that the quadratic
approximation improves near a minimum and gives a conjugate gradient method fast
convergence. On Sinc, \gls{scg} attained the target with less work than \gls{sgd}.

The order reversed in wall-clock time, with \gls{sgd} the slowest algorithm on
every problem by a factor between 1.4 on Sin2D and 24.7 on Cancer. An epoch of \gls{sgd}
performs many small matrix products where a batch method performs one large
product, and each product carries a fixed overhead that the \glspl{cu} ignore.

\subsection{Overall Comparison}

Table~\ref{tab:friedman} reports the tests across the nine problems and the mean
ranks. Bold marks the lowest mean rank of a criterion, and a dagger marks a mean
rank that exceeds the lowest by more than the critical difference of 1.10. The work to target and the wall-clock time differed significantly, whereas the
generalisation error and the \gls{mse} on the test set did not. On
generalisation error, \gls{scg} and \gls{sgd} shared a mean rank of 1.78. Within the problems, \gls{scg} won three significant comparisons and lost one,
LeapFrog won three and lost two, and \gls{sgd} won one and lost four. On work to target,
\gls{sgd} was significantly faster than LeapFrog, while the lead of \gls{sgd} over
\gls{scg}, at 1.00 mean rank, fell just within the critical difference of 1.10. On
wall-clock time, both batch methods were significantly faster than \gls{sgd}.

\begin{table}[t]
\centering
\caption{Tests across the nine problems and mean ranks}
\label{tab:friedman}
\footnotesize
\setlength{\tabcolsep}{3pt}
\begin{tabular}{@{}l r rrr@{}}
\toprule
& & \multicolumn{3}{c}{Mean rank} \\\cmidrule(l){3-5}Criterion & \shortstack{Probability\\value} & SGD & SCG & LeapFrog \\
\midrule
Generalisation & 0.2636 & \textbf{1.78} & \textbf{1.78} & 2.44 \\
Squared error & 0.1690 & \textbf{1.56} & 2.00 & 2.44 \\
Work to target & 0.0018 & \textbf{1.11} & 2.11 & 2.78$^\dagger$ \\
Wall-clock time & 0.0009 & 3.00$^\dagger$ & 1.67 & \textbf{1.33} \\
\bottomrule
\end{tabular}

\end{table}

The best algorithm therefore depends on the scarce resource. \Gls{sgd} is the best
choice when gradient evaluations are scarce and tuning is affordable, whereas the
batch methods are the better choice when time is scarce. Between the two batch
methods, the effort of tuning decides. The default parameters of \gls{scg} were statistically as
good as the best setting on every problem, and \gls{scg} did not fail on Digits.
\Gls{scg} was therefore the most practical of the three algorithms. The finding that
\gls{sgd} needed less work than \gls{scg} does not contradict
M{\o}ller~\cite{moller1993}, who found \gls{scg} to need between 8.4
and 23.4 times fewer \glspl{cu} than backpropagation. The backpropagation of
M{\o}ller updated the weights once per pass over the data, at a learning rate set
by hand and a momentum of 0.9. The \gls{sgd} of the present study used small
batches and was tuned over 45 settings, a far stronger baseline.

The repetition with \glspl{relu} confirmed that no algorithm generalised
significantly better, at a Friedman probability value of 0.46, and the ranks on
wall-clock time were identical to the ranks of the sigmoid study. The advantage of \gls{sgd} in work to target lost significance, with a
probability value of 0.24, so the fast start of \gls{sgd} depended on the
activation function. The \gls{relu} networks required more hidden units on seven
problems, up to 96 on Sinc and Sin2D, because a piecewise linear network
approximates a smooth curve with many linear pieces. The \glspl{relu} halved the
error on Digits for all three algorithms, whereas the sigmoid units were better in five of the nine algorithm--problem
pairs of the smooth problems Sinc, Sin2D and Friedman1.

In summary, the optimal number of hidden units ranged from one to 24. \Gls{scg} was free of parameters in practice, whereas LeapFrog suffered from one
interaction of parameters and \gls{sgd} was sensitive to the learning rate. No algorithm
generalised significantly better across the nine problems. \Gls{sgd} was the
fastest per unit of work, and the batch methods were the fastest in wall-clock
time.

\glsresetall

\section{Conclusion}

The study compared \gls{sgd}, \gls{scg} and LeapFrog for the training of feedforward
neural networks with one hidden layer, on five classification and four function
approximation problems. After the number of hidden units and the control
parameters had been chosen on validation data, the algorithms were compared on 30
fresh paired runs per problem under an equal budget of \glspl{cu}.

No algorithm generalised significantly better across the nine problems.
\Gls{sgd} attained a validation cost common to all three algorithms with
significantly less work, but required
significantly more wall-clock time. \Gls{scg} was insensitive to the control parameters over the grid tested, and
the default parameters of \gls{scg} were statistically as good as the best
setting on every problem. The tuned learning rate of \gls{sgd} lay close
to the boundary of stability and failed on larger networks. LeapFrog failed to
optimise the large Digits network because the step limit dominated the
trajectory. \Gls{scg} was concluded to be the most
practical algorithm for the problems studied, and the absence of a significant
difference in generalisation held for \glspl{relu} as well.

Future work could address four questions. A LeapFrog step limit that scales with
the number of weights could remove the failure on large networks. A comparison over a
range of budgets would show where the fast start of \gls{sgd} stops mattering. A
learning rate that adapts to the curvature would strengthen the \gls{sgd} baseline.
A cross-entropy cost function would test whether the conclusions depend on the cost
function.

\bibliographystyle{IEEEtranS}
\bibliography{references}

\end{document}
